\documentclass[11pt]{article}
\usepackage[T1]{fontenc}
\usepackage{amsmath,amssymb,amsthm}
\usepackage[margin=1in]{geometry}
\usepackage[hidelinks]{hyperref}
\newtheorem{theorem}{Theorem}
\newtheorem{lemma}[theorem]{Lemma}
\newtheorem{proposition}[theorem]{Proposition}
\newtheorem{corollary}[theorem]{Corollary}
\newcommand{\C}{\mathcal C}
\newcommand{\D}{\mathcal D}
\newcommand{\aideal}{\mathfrak a}
\newcommand{\Nideal}{\mathcal N_G}
\newcommand{\R}{\mathbb R}
\title{Finite-group quotients and geometric obstacle envelopes\\
Full semivaluation fibres, normal singularities, and canonical norm bounds}
\author{Independent scoped ORIGINAL P12.3 deductions; unreviewed draft}
\date{9 October 2026}
\begin{document}
\maketitle
\begin{abstract}
Relative to the entire qualified geometric comparison G331, we prove
that the full normalized boundary of a finite invariant-ring extension
is its finite-group orbit quotient, including every nonminimal support.
The proof uses coefficient ideals of an orbit product and compact lattice
density, rather than an unproved transitivity theorem for valuation
extensions. Finite orbit maxima show that raw geometric envelopes commute
with pullback. Greatest continuous or decreasing-completed minorants exist
exactly when the corresponding invariant source obstacle has one.
Consequently all finite quotients of smooth affine closed germs over an
algebraically closed field of characteristic zero have continuous geometric
envelopes, in every dimension, including nonisolated normal singularities.
Coefficient norms also sandwich the every-degree canonical ideal profiles
and give a quantitative all-degree convergence bound. General excellent
normal higher-dimensional regularity and arbitrary conductor-fibre
classification remain open. This is an independently written scoped draft,
not a closure of the unchanged original problem.
\end{abstract}

\section{Exact domain and retained geometric meaning}
Let $R$ be an excellent normal local domain of positive dimension, and
let $S$ be an excellent normal semilocal domain, finite over $R$.
All its marked closed centres are the centres above the maximal ideal
of $R$. A finite group $G$ acts on $S$ and $R=S^G$; no faithfulness
or invertibility of $h=|G|$ in the ring is required. Fix a proper primary
ideal $\aideal\subset R$. Set
\[
 V_R=V_{\aideal}(R),\qquad V_S=V_{\aideal S}(S),\qquad
 \rho:V_S\longrightarrow V_R.
\]
The source is the finite disjoint union of its local-centre sections.
The full pointwise compact topology, every prime support, and infinite
values of individual functions are retained. Restriction is surjective
and is a compact Hausdorff quotient by the full qualified G331 boundary
theorem. In particular no semivaluation with nonzero support is extended
to the original fraction field.

The classes $\C_R,\D_R$ and their source versions are the independently
geometric classes of G331. Their core consists of actual negative
effective finite real Cartier presentations on directed actual primary
blowups, with degree at least zero on every contracted integral curve.
Only after constructing that core does G331 prove comparison with finite
negative ideal profiles and their uniform and actual decreasing-net
completions. We use that proved comparison, not a new positivity
definition. A proper primary profile $-v(I)/k$ has an actual negative
effective relatively ample presentation before any limit; unit ideals
give zero. The source product classes agree with semilocal ideal profiles
for ideals containing a power of $\aideal S$.

The sole incoming correctness grant is ENTIRE qualified G331, including
all five fields and its mandatory source, model, support, historical-audit
and radial qualifications. Its exact boundary, comparison, finite
positive descent, gauge and completion statements are used. Previous
envelope papers are retained proof DATA and their exact applications are
rejudged. General model cofinality, singular higher-dimensional resolution,
finite generation and radial CLD identification are not inferred.

\section{Orbit coefficient ideals on every support}
For an $S$-ideal $J$ containing a power of $\aideal S$, choose finite
$R$-module generators $j_1,\ldots,j_l$ of $J$; these exist because
$S$ is finite over the Noetherian ring $R$. Define $\Nideal(J)$ to
be the coefficient ideal in $R$ of
\[
 P_J(T)=\prod_{\gamma\in G}
              \left(\sum_{i=1}^l T_i\gamma(j_i)\right).
\]
Every coefficient is invariant because $G$ permutes the factors, hence
belongs to $R$. Changing $R$-module generators amounts to a linear
substitution with coefficients in $R$. Such coefficients are fixed by
every group element. Substitution in both directions proves independence
of the coefficient ideal from the generators.

\begin{lemma}\label{orbit:norm}
The ideal $\Nideal(J)$ is unit-or-primary and, for every $w\in V_S$,
\[
 \rho(w)(\Nideal(J))=\sum_{\gamma\in G}(\gamma^*w)(J),
 \qquad (\gamma^*w)(s)=w(\gamma(s)).
\]
The identity includes all source supports and involves only finite ideal
values. It requires no defectless ramification formula.
\end{lemma}
\begin{proof}
If $J\supset\aideal^N S$, add the finitely many base generators
$b$ of $\aideal^N$ to its $R$-module generators. Their pure variable
coefficients are $b^h$, so the coefficient ideal contains an ideal with
radical the maximal ideal of $R$. Thus it is unit-or-primary. If $J=S$,
include $1$ and get the unit ideal. If $J$ is proper at some marked
centre $M$, every coefficient of $P_J$ lies in $M$, because the identity
factor vanishes modulo $M$; its invariant coefficients lie in the maximal
ideal of $R$. This also identifies the proper case.

Quotient $S$ by the support of $w$ and enter the resulting valued
quotient fraction field. The minimum of the coefficient values of a
polynomial is multiplicative: divide each nonzero polynomial by a
least-valued coefficient and reduce to nonzero polynomials over the
residue field, whose product is nonzero. A polynomial zero in the
quotient has infinite coefficient minimum. Here no factor is zero:
$\gamma^*w$ has the invariant gauge $\aideal S$ of value one, and
$J$ contains its power. The minimum for each factor is exactly
$(\gamma^*w)(J)$, since the chosen elements generate the entire ideal
even as an $R$-module. Taking the coefficient minimum of the product
gives the formula. This proves the statement without treating a
semivaluation as a valuation of $\operatorname{Frac}(S)$.
\end{proof}

\begin{lemma}\label{orbit:fibres}
Every full restriction fibre is exactly one $G$-orbit. Hence
$V_S/G\to V_R$ is a homeomorphism.
\end{lemma}
\begin{proof}
For $w\in V_S$ put
\[
 \mu_w=\frac1h\sum_{\gamma\in G}\delta_{\gamma^*w}.
\]
If $\rho(w)=\rho(w')$, Lemma \ref{orbit:norm} shows that these two
probability measures have equal integrals against every boundary-power
ideal gauge. The rational difference span of such gauges is a lattice:
after a common denominator, its minima use
\[
 \min\{v(I)-v(J),v(K)-v(H)\}=v(IH+KJ)-v(JH).
\]
It contains constants from $v(\aideal S)=1$, separates all coordinates
by $v((f)+(\aideal S)^N)=\min(v(f),N)$, and separates distinct
semilocal centres by centre-primary ideals. Its uniform closure is all
$C(V_S,\R)$ by the compact finite maximum/minimum interpolation
argument of G331. Thus the two finite orbit averages agree on every
continuous function. Distinct finite subsets of a compact Hausdorff space
can be separated by a continuous function; all orbit weights are positive.
Their supports therefore coincide. The two orbits are the same.
Conversely every group translate restricts identically to $R$. The
surjective compact quotient then gives the asserted homeomorphism.
No orbit statement is restricted to minimal primes or divisorial points.
\end{proof}

The averaging operator $T$ defined by $\rho^*T(u)=h^{-1}
\sum_{\gamma}\gamma^*u$ has positive norm one on continuous functions.
On negative rational profiles Lemma \ref{orbit:norm} gives
\[
 T(-w(J)/k)=-v(\Nideal(J))/(hk).
\]
Uniform comparison proves $T(\C_S)\subset\C_R$. For an actual
decreasing $\C_S$-net, a finite positive sum commutes with its extended
nonpositive limit: a bottom summand forces bottom, and otherwise a common
late index approximates all finitely many summands. Hence
$T(\D_S)\subset\D_R$ with an actual decreasing presentation. These
are positive cone statements, not merely descent of continuous values.

\section{Exact envelope transport}
For a finite continuous signed obstacle $f$ on $V_R$ define separately
the raw suprema over its $\C$ and $\D$ minorants. They are finite,
since a sufficiently negative constant is a minorant and zero bounds
the cone above. Denote their common value, established below, by $p_R(f)$.

\begin{theorem}\label{orbit:envelope}
For every such $f$, on every point of $V_S$,
\[
 \rho^*p_R(f)=p_S(\rho^*f).
\]
The equality holds for each of the two raw classes independently.
A greatest $\C$ or $\D$ minorant downstairs exists if and only if
the pulled-back invariant source obstacle has one. When either exists,
both greatest minorants are the same continuous geometric function.
Consequently source universal continuous-envelope regularity implies
target universal continuous-envelope regularity. The converse requires
only source invariant obstacles, not all source obstacles.
\end{theorem}
\begin{proof}
Pullback preserves both classes and gives one inequality. If $u$ is any
source minorant of $\rho^*f$, its finite orbit maximum
\[
 u^G=\max_{\gamma\in G}\gamma^*u
\]
belongs to the same class, is invariant, remains below the obstacle,
and dominates $u$. It descends positively, either by the full-fibre
G331 theorem or by $T(u^G)$ above, whose pullback is $u^G$.
It therefore bounds $u$ by the pulled-back target raw supremum.
Taking the supremum proves the equality. No increasing-limit positivity
or fibrewise-minimum positivity has been assumed.

If the source has a greatest minorant, uniqueness and invariance of the
obstacle make it invariant; positive descent gives the target greatest
minorant. If the target has one, its pullback equals the source raw
supremum by the proved formula, and so is the source greatest minorant.
The continuity and equality of the two classes follow from the next
canonical argument. It applies to each ring independently.
\end{proof}

\section{Canonical profiles and quantitative norm comparison}
For either ring and its continuous signed obstacle put
\[
 A_m(f)=\{s:v(s)+mf(v)\ge0\text{ for every normalized }v\},
 \qquad h_m(f)=-v(A_m(f))/m,\quad m\ge1.
\]
The ultrametric inequality makes this an ideal. It contains a sufficiently
high boundary power, so is finitely generated and unit-or-primary.
Every $h_m$ is an actual geometric continuous minorant of $f$, with
ample primary presentation in the proper case, and
$A_mA_n\subset A_{m+n}$.

Here is the full raw-supremum argument, rather than an appeal to finite
generation. If $g\in\C$ and $g\le f$, choose a rational primary
profile $q=-v(I)/k$ with $\|q-g\|<\epsilon$. Then
$I^N\aideal^{\lceil Nk\epsilon\rceil}\subset A_{Nk}(f)$,
so $h_{Nk}\ge q-\lceil Nk\epsilon\rceil/(Nk)$ and
$\sup_m h_m\ge g$. If $g=\inf_i g_i\in\D$ and $g\le f$,
the open sets $\{g_i<f+\epsilon\}$ cover the compact section.
A common index from a finite subcover has $g_i-\epsilon\le f$ and
is a $\C$ minorant dominating $g-\epsilon$. Thus
\[
 p_C(f)=p_D(f)=\sup_m h_m(f).
\]
Write $M_1=\|v(A_1(f))\|_\infty$. For fixed $m$ and
$n=km+r$, $0\le r<m$, the graded inclusion gives
$h_n\ge(km/n)h_m-rM_1/n$. Hence every degree tends pointwise to
the raw supremum. If that supremum is continuous, finitely many strict
open tests $h_m>p-\epsilon$ cover the section; a common multiple $N$
of their indices has $h_N>p-\epsilon$ everywhere. The same graded
bound gives $0\le p-h_n\le\epsilon+NM_1/n$ for late $n$.
Thus convergence is uniform. Conversely a uniform limit is continuous
and in $\C$. A greatest $\D$ minorant equals the raw supremum,
which is both lower semicontinuous as a supremum of continuous functions
and upper semicontinuous as a decreasing completion; it is continuous.
This completes the existence claims in Theorem \ref{orbit:envelope}.

Now write $J_m=A_m^S(\rho^*f)$ and $I_m=A_m^R(f)$.
Surjectivity of the full boundary gives the exact contraction
$I_m=J_m\cap R$; $J_m$ is $G$-stable. The coefficient norm yields
\[
 \Nideal(J_m)\subset I_{hm},\qquad
 \rho^*h_m^R\le h_m^S\le\rho^*h_{hm}^R.
\]
Indeed $I_mS\subset J_m$ proves the first profile inequality.
For the second, every coefficient of $P_{J_m}$ is a sum of products
of $h$ elements of $J_m$, hence lies in $J_m^h\cap R\subset I_{hm}$.
Stability and Lemma \ref{orbit:norm} give
$w(\Nideal(J_m))=h\,w(J_m)$.

\begin{proposition}\label{orbit:error}
Suppose $p_S(\rho^*f)$ is continuous and put
$\beta_m=\|p_S-h_m^S\|_\infty$. For $n\ge h$, set
$m=\lfloor n/h\rfloor$. Then
\[
 \|p_R-h_n^R\|_\infty
 \le \beta_m+\frac{h-1}{n}\|v(I_1)\|_\infty.
\]
Every degree converges uniformly. No finite generation is required.
\end{proposition}
\begin{proof}
Write $n=hm+r$, $0\le r<h$. The actual ideal
$\Nideal(J_m)I_1^r$ is contained in $I_n$, so
\[
 \rho^*h_n^R\ge \frac{hm}{n}h_m^S-\frac{r}{n}w(I_1).
\]
Subtract from $\rho^*p_R=p_S\le0$ to obtain the bound. The source
$\beta_m$ tends to zero by the preceding canonical proof.
Surjectivity makes pullback an isometry, so the source supremum norm
is the target supremum norm. All ideals and geometric primary
presentations are actual before the convergence argument.
\end{proof}

\section{Normal quotient singularities in every dimension}
Let $k$ be algebraically closed of characteristic zero, and let a finite
group act by $k$-automorphisms on a smooth affine $k$-variety, fixing a
closed smooth point. Localize the invariant affine ring at its image.
Its finite source has just the marked source centre. The source belongs
to the exact published smooth affine CLOSED BFJ domain. The qualified
smooth-envelope proof is recalled below. Theorem \ref{orbit:envelope}
therefore proves universal continuous geometric-envelope regularity on
the normal quotient germ in every dimension.

For clarity the ring assertions do not rely on a quoted invariant-theory
label. If $B$ is a finitely generated coordinate domain and $G$ acts,
each algebra generator satisfies its monic orbit polynomial.
Let $B_0$ be the $k$-algebra generated by all its coefficients.
Then $B$ is finite over the Noetherian ring $B_0$.
The intermediate invariant ring $B^G$ is a finite $B_0$-module, hence
finitely generated over $k$, and $B$ is finite over it.
For an invariant fraction $a/b$, the invariant denominator
$\prod_\gamma\gamma(b)$ clears it to an invariant element of $B$.
Thus $\operatorname{Frac}(B)^G=\operatorname{Frac}(B^G)$.
If $B$ is normal, an element of this fraction field integral over
$B^G$ is integral over $B$, belongs to $B$, and is invariant.
Hence the invariant ring is normal; finite-type local rings over a field
are excellent.

If a second maximal ideal above the invariant image were not in the
orbit of the fixed point, Chinese remainder gives an element zero at
the fixed point and one at every translate of that second ideal.
Its orbit product is invariant, vanishes at the fixed point and is one
at the second ideal, contradicting equal contraction. Thus that fibre
has just the fixed point. Inverting invariant denominators gives the
local source; any further local denominator is a unit because this is
the only maximal ideal. The same orbit-product clearing proves that
local invariants equal the localized invariant ring. Dimension is
preserved by the finite integral extension.

For the standard example take $B=k[x_1,\ldots,x_d]$ and let the
order-two action negate $x_1,x_2$ and fix the other coordinates.
For $d\ge3$ the invariant local germ is
\[
 R=k[a,b,c,x_3,\ldots,x_d]_0/(ac-b^2),\qquad
 (a,b,c)=(x_1^2,x_1x_2,x_2^2).
\]
It is a normal nonisolated singularity and has continuous geometric
envelopes for ALL continuous signed obstacles. Normality follows from
the invariant-ring proof. The displayed generators follow by grouping
monomials with even total degree in $x_1,x_2$; their sole relation is
$ac=b^2$, seen by the two parity normal forms. The singular locus is
the axis $a=b=c=0$, of dimension $d-2$; the Jacobian or tangent
dimension detects singularity there. This is genuinely a NORMAL
higher-dimensional sector, not another conductor counterexample.

\section{Exact source application, completion, and remaining original scope}
The smooth source regularity used here is at the complete previously
checked BFJ setting, simple-modification domain, closed-simplex and
all-retraction scope. For a signed primary-model obstacle on the smooth
source, its function is determined on a common simple dual complex.
BFJ 3.15--3.18 constructs the nef envelope there. The WHOLE 3.18 proof
tests multiplier inequalities on that fixed complex, where thinness is
bounded; it does not claim thinness bounded on the full boundary.
Use its real parameter $t=m+C$, not an unjustified integer replacement.
Its estimate gives the bounded envelope between canonical ideal
profiles with uniform error at most $C\|p\|_\infty/m$.
Closed-simplex continuity and dense divisorial values extend the
inequalities to every face; taking all retraction infima extends them
to EVERY semivaluation, including nonminimal support. Completed primary
ideals contract exactly. The zero case uses the unit ideal. Uniform
signed primary-model lattice approximation of an arbitrary continuous
obstacle and the sup-norm contraction of raw envelopes give the general
continuous source result. Historical smooth principalization, multiplier
subadditivity and Izumi remain disclosed published background, not a
fresh full recursive audit.

This application retains ENTIRE qualified G331 and the complete current
normal-surface/smooth proof carrier. The complete BFJ setting,
simple-modification domain and 3.15--3.25 immediate chain have been freshly
physically reread for this application; other complete immediate units
and formula pixels retain their earlier qualifications without fresh
credit. No primary literal body or formula pixel is embedded here.
The new coefficient-product, orbit-measure, envelope and norm-filtration
proofs are given above; no new Galois valuation-transitivity source is
needed. Normal source CE is an antecedent of the abstract theorem and
is supplied only in the stated proved sectors.

For any excellent quotient germ just treated, G331's full completion
homeomorphism and exact primary contraction carry the entire obstacle,
cone and envelope result to its formal completion. This does not assert
equivalence of geometric model categories or identify an arbitrary
formal invariant ring with such a completion without a further proof.
In dimension two the abstract finite-group result also applies whenever
the source normal-surface CE theorem is used at its retained
arbitrary-field/equal-or-mixed-characteristic qualifications; this adds
no higher-dimensional positive-characteristic source regularity.

For another target proper primary gauge $\mathfrak b$, divide all
homogeneous functions by the same bounded positive denominator
$v(\mathfrak b)$. It is constant on every restriction fibre and
$G$-orbit. The envelope equality, continuity, positive descent and
greatest-minorant criterion transport exactly. The canonical ideals
remain unchanged for the transported obstacle, since their inequalities
are divided by this same positive number. Zero-dimensional sections
are empty and have the unique empty-function statement directly.

General excellent NORMAL dimension at least three CE, positive- or
mixed-characteristic higher-dimensional regular-source CE, and
necessary-and-sufficient envelope criteria for arbitrary nonnormal
conductor fibres remain OPEN. A general singular germ is not assumed
to admit a finite regular quotient cover. Products, energy, full finite
operators and literal radial CLD comparison remain separate targets.
This paper was developed independently after the native296297 launches
without reading their mathematics or sending hot feedback. Compilation
and acceptance are unverified. The unchanged original P12.3 remains OPEN;
candidate is false.
\end{document}
